\begin{Solbook}{1}
 \par \noindent  $x=-1$  \hfill \qedEX 
\end{Solbook}
\begin{Solbook}{2}
 \par \noindent  a) $A=-\dfrac {3}{5}$ \hspace {2cm} b) $B=\dfrac {\sqrt {x}}{\sqrt {x}-4}$ \hspace {2cm} c) $0<x\le 4$  \hfill \qedEX 
\end{Solbook}
\begin{Solbook}{3}
 \par \noindent Số học sinh dự thi của trường $A$ là 370. Số học sinh dự thi của trường $B$ là 310. \hfill \qedEX 
\end{Solbook}
\begin{Solbook}{4}
 \par \noindent Độ cao của Flycam là khoảng 466,7 m. \hfill \qedEX 
\end{Solbook}
\begin{Solbook}{5}
 \par \noindent \immini {\begin {enumerate} \item $B,E,C,D$ cùng thuộc đường tròn $(I)$. \item $CB$ là đường kính, $DE$ là dây cung khác đường kính. \par Suy ra $CB>ED$. \item $BKCH$ là hình bình hành suy ra $CK \parallel BD$. \item $\dfrac {\pi \cdot 3\cdot 70}{360}-\dfrac {1}{2}\sqrt {3}\cdot \sin 55^\circ \cdot 2\cdot \sqrt {3}\cdot \cos 55^\circ \approx 0,4\, \text {cm}^2$ . \end {enumerate}} {\begin {tikzpicture}[scale=1, font=\scriptsize ] \def \R {2.5} \coordinate (B) at (-2.5,0); \coordinate (C) at (2.5,0); \coordinate (I) at (0,0); \coordinate (A) at (-0.5,4); \draw (I) circle(\R ); \path [name path=cirI] (I) circle(\R ); \path [name path=ab] (A)--(B); \path [name intersections={of=cirI and ab, by={E1,E2}}]; \coordinate (E) at (E1); \path [name path=ac] (A)--(C); \path [name intersections={of=cirI and ac, by={D1,D2}}]; \coordinate (D) at (D1); \path [name path=bd] (B)--(D); \path [name path=ce] (C)--(E); \path [name intersections={of=bd and ce, by=H}]; \draw (A)--(B)--(C)--cycle; \draw (C)--(E)--(D)--(B); \coordinate (K) at ($-1*(H)$); \draw (B)--(K)--(C); \draw (H)--(K); \draw (I)--(D); \draw (H)--(A); \foreach \p /\pos in {A/above,B/left,I/below,C/right,E/above left,D/above right,H/right,K/below} {\filldraw [fill=white] (\p ) circle(1.5pt) node[\pos ] {$\p $}; } \end {tikzpicture}}  \hfill \qedEX 
\end{Solbook}
\begin{Solbook}{6}
 \par \noindent  Gợi ý: \par Chứng minh $2(a^2+b^2+c^2)-2(ac+ac+bc)\ge 0$ \par Dùng bất đằng thức trong tam giác $a+b>c,a+c>b,b+c>a$ bình phương hai vế của các bất đẳng thức trên, cộng chúng lại và biến đổi.  \hfill \qedEX 
\end{Solbook}
